{\scriptsize\tiny
    \begin{minipage}[t]{0.54\textwidth}
        \begin{tabularx}{\linewidth}{|p{0.35\textwidth}|X|}
            \hline 
            \textbf{Notions et contenus} & \textbf{Capacités exigibles} \\
            \hline 
            \multicolumn{2}{|l|}{5.2.2. Exemples de champs électrostatiques } \\
            \hline 
            Dipôle électrostatique. Moment dipolaire.  & 
            Citer les conditions de l'approximation dipolaire. \\
            \hline 
            Potentiel et champ créés par un dipôle. & Établir l'expression du potentiel électrostatique.
            Comparer la décroissance du champ et du potentiel avec la distance dans le cas d'une charge ponctuelle et dans le cas d'un dipôle.
            Tracer l'allure des lignes de champ électrostatique engendrées par un dipôle. \\
            \hline 
            Actions subies par un dipôle placé dans un champ électrostatique d'origine extérieure : résultante et moment. & Utiliser les expressions fournies de la résultante et du moment des actions subies par un dipôle placé dans un champ électrostatique d'origine extérieure.\\
            \hline
            Énergie potentielle d'un dipôle rigide dans un champ électrostatique d'origine extérieure. & Utiliser l'expression fournie de l'énergie potentielle d'un dipôle rigide dans un champ électrostatique d'origine extérieure. Prévoir qualitativement l'évolution d'un dipole rigide dans un champ électrostatique d'origine extérieure. \\
        \hline 
        \end{tabularx}
\end{minipage}
\begin{minipage}[t]{0.43\textwidth}
    \begin{tabularx}{\linewidth}{|p{0.35\textwidth}|X|}
        \hline 
        
        Interactions ion-molécule et molécule-molécule. & Expliquer qualitativement la solvatation des ions dans un solvant polaire. \\
        \hline 
        Dipôle induit. Polarisabilité. & Associer la polarisabilité et le volume de l'atome en ordre de grandeur. \\ 
        \hline
        Plan infini uniformément chargé en surface. & Établir l'expression du champ créé par un plan infini uniformément chargé en surface.\\ 
        \hline
        Établir l'expression du champ créé par un plan infini uniformément chargé en surface. & Établir l'expression du champ créé par un condensateur plan.
        Déterminer l'expression de la capacité d'un condensateur plan.
        Citer l'ordre de grandeur du champ disruptif dans l'air.
        Déterminer l'expression de la densité volumique d'énergie électrostatique dans le cas du condensateur plan à partir de celle de l'énergie du condensateur. \\ 
        \hline
        Énergie de constitution d'un noyau atomique modélisé par une boule uniformément chargée. & Énergie de constitution d'un noyau atomique modélisé par une boule uniformément chargée. \\ 
        \hline
    \end{tabularx}
\end{minipage}
}
